% =============================================================================
% SECCION: VALIDACION ESTADISTICA
% =============================================================================

\section{Validaci\'on Estad\'istica}
\label{sec:statistical_validation}

La evaluaci\'on del rendimiento de estrategias de trading requiere rigor estad\'istico m\'as all\'a de simples comparaciones de retornos. Esta secci\'on presenta los tests estad\'isticos aplicados para validar la significancia de los resultados.

\subsection{Bootstrap de Sharpe Ratios}
\label{subsec:bootstrap_sharpe}

\subsubsection{Metodolog\'ia}

El Sharpe Ratio est\'andar no provee informaci\'on sobre la precisi\'on de la estimaci\'on. Utilizamos bootstrap no param\'etrico \citep{efron1979bootstrap} para construir intervalos de confianza:

\begin{enumerate}
    \item Generar $B = 10,000$ muestras bootstrap con reemplazo de los retornos diarios
    \item Calcular el Sharpe Ratio para cada muestra bootstrap: $\hat{S}_b, \; b = 1, \ldots, B$
    \item Construir el intervalo de confianza al 95\% usando percentiles: $[\hat{S}_{0.025}, \hat{S}_{0.975}]$
\end{enumerate}

\subsubsection{Correcci\'on de Ledoit-Wolf}

Aplicamos la correcci\'on de \cite{ledoit2008robust} para muestras finitas y autocorrelaci\'on en los retornos:

\begin{equation}
\text{SE}(\hat{S}) = \sqrt{\frac{1 + \frac{\hat{S}^2}{2} \cdot (1 + \hat{\rho}_1)}{T \cdot (1 - \hat{\rho}_1)}}
\end{equation}

donde $\hat{\rho}_1$ es la autocorrelaci\'on de primer orden de los retornos y $T$ el n\'umero de observaciones.

\subsubsection{Resultados: Intervalos de Confianza}

\begin{table}[htbp]
\centering
\caption{Intervalos de Confianza del Sharpe Ratio (Bootstrap 95\%)}
\label{tab:sharpe_ci}
\small
\begin{tabular}{@{}lrrrrr@{}}
\toprule
\textbf{Modelo} & \textbf{Sharpe} & \textbf{IC 95\% Inf} & \textbf{IC 95\% Sup} & \textbf{P(S$>$0)} & \textbf{P(S$>$SPY)} \\
\midrule
Ridge & 0.932 & 0.612 & 1.248 & 99.8\% & 58.2\% \\
CNN\_LSTM & 0.744 & 0.428 & 1.056 & 99.1\% & 42.6\% \\
BiGRU & 0.680 & 0.374 & 0.982 & 98.4\% & 36.1\% \\
GradientBoosting & 0.589 & 0.214 & 0.960 & 96.8\% & 28.3\% \\
N-BEATS & 0.511 & 0.198 & 0.820 & 94.7\% & 18.9\% \\
\midrule
SPY (B\&H) & 0.882 & 0.584 & 1.176 & 99.6\% & -- \\
\bottomrule
\end{tabular}
\vspace{0.2cm}
\small\textit{
\small
Nota: P(S$>$0) = probabilidad de Sharpe positivo. P(S$>$SPY) = probabilidad de superar el Sharpe del benchmark. IC calculados con 10,000 iteraciones bootstrap.}
\end{table}

\subsubsection{Interpretaci\'on}

Los intervalos de confianza revelan que:

\begin{itemize}
    \item \textbf{Ridge}: Con 99.8\% de probabilidad, tiene Sharpe positivo. Sin embargo, solo 58.2\% de probabilidad de superar al benchmark SPY.
    \item \textbf{CNN-LSTM}: Sharpe positivo con 99.1\% de confianza, pero solo 42.6\% de superar SPY.
    \item \textbf{Solapamiento significativo}: Los intervalos de confianza de los mejores modelos se solapan sustancialmente con el del benchmark, indicando que la superioridad no es estad\'isticamente conclusiva.
\end{itemize}

\subsection{Test de Diebold-Mariano}
\label{subsec:diebold_mariano}

\subsubsection{Formulaci\'on}

El test de \cite{diebold1995comparing} eval\'ua si la diferencia en error de predicci\'on entre dos modelos es estad\'isticamente significativa:

\begin{equation}
DM = \frac{\bar{d}}{\sqrt{\hat{V}(\bar{d})}}
\end{equation}

donde $d_t = L(e_{1,t}) - L(e_{2,t})$ es la diferencia en p\'erdida entre los dos modelos en el tiempo $t$, y $\hat{V}(\bar{d})$ es un estimador consistente de la varianza de $\bar{d}$ que considera autocorrelaci\'on.

Bajo la hip\'otesis nula de igual capacidad predictiva, $DM \xrightarrow{d} N(0,1)$.

\subsubsection{Resultados: Comparaci\'on con Benchmark}

\begin{table}[htbp]
\centering
\caption{Test de Diebold-Mariano vs Modelo Naive (Retorno = 0)}
\label{tab:dm_test}
\small
\begin{tabular}{@{}lrrr@{}}
\toprule
\textbf{Modelo} & \textbf{Estad\'istico DM} & \textbf{p-valor} & \textbf{Significativo (5\%)} \\
\midrule
Ridge & 2.847 & 0.0044 & S\'i \\
GradientBoosting & 1.923 & 0.0545 & No \\
CNN\_LSTM & 2.412 & 0.0159 & S\'i \\
BiGRU & 2.156 & 0.0311 & S\'i \\
RandomForest & 1.284 & 0.1992 & No \\
\bottomrule
\end{tabular}
\vspace{0.2cm}
\small\textit{
\small
Nota: Estad\'istico DM positivo indica que el modelo supera al benchmark naive. Correcci\'on Newey-West con 5 rezagos para autocorrelaci\'on.}
\end{table}

\subsection{Model Confidence Set (MCS)}
\label{subsec:mcs}

\subsubsection{Metodolog\'ia}

El Model Confidence Set de \cite{hansen2011model} identifica el conjunto de modelos que contiene al mejor modelo verdadero con una probabilidad dada:

\begin{enumerate}
    \item Comenzar con el conjunto completo de modelos $\mathcal{M}_0$
    \item Testear la hip\'otesis de igual capacidad predictiva (EPA) sobre $\mathcal{M}_0$
    \item Si EPA se rechaza, eliminar el peor modelo
    \item Repetir hasta que EPA no se rechace
\end{enumerate}

El MCS final $\hat{\mathcal{M}}^*_{\alpha}$ contiene todos los modelos que no pueden ser descartados como inferiores al nivel $\alpha$.

\subsubsection{Resultados}

Al nivel de confianza del 90\%, el Model Confidence Set contiene:

\begin{equation}
\hat{\mathcal{M}}^*_{0.10} = \{\text{Ridge}, \text{CNN-LSTM}, \text{BiGRU}, \text{GradientBoosting}\}
\end{equation}

Esto indica que estos cuatro modelos no pueden distinguirse estad\'isticamente como superiores entre s\'i, pero s\'i son estad\'isticamente superiores al resto.

\subsection{Test de Clark-West}
\label{subsec:clark_west}

\subsubsection{Formulaci\'on}

El test de \cite{clark2007approximately} es apropiado para comparar modelos anidados (nested models):

\begin{equation}
CW_t = (e_{1,t}^2 - e_{2,t}^2) + (e_{1,t} - e_{2,t})^2
\end{equation}

donde el t\'ermino de ajuste $(e_{1,t} - e_{2,t})^2$ corrige el sesgo que surge cuando el modelo restringido es verdadero.

\subsubsection{Aplicaci\'on}

Comparamos cada modelo contra el modelo naive (predicci\'on = promedio hist\'orico):

\begin{table}[htbp]
\centering
\caption{Test de Clark-West vs Modelo de Media Hist\'orica}
\label{tab:cw_test}
\small
\begin{tabular}{@{}lrrr@{}}
\toprule
\textbf{Modelo} & \textbf{CW Stat} & \textbf{p-valor (1-cola)} & \textbf{Conclusi\'on} \\
\midrule
Ridge & 2.31 & 0.010 & Rechaza H$_0$ \\
CNN\_LSTM & 1.89 & 0.029 & Rechaza H$_0$ \\
BiGRU & 1.67 & 0.047 & Rechaza H$_0$ \\
GradientBoosting & 1.52 & 0.064 & Marginal \\
\bottomrule
\end{tabular}
\vspace{0.2cm}
\small\textit{
\small
Nota: H$_0$: el modelo extendido no tiene capacidad predictiva superior al modelo restringido. Test de una cola dado que la hip\'otesis alternativa es unidireccional.}
\end{table}

\subsection{An\'alisis de Overfitting}
\label{subsec:overfitting}

\subsubsection{Comparaci\'on In-Sample vs Out-of-Sample}

Una se\~nal de overfitting es una degradaci\'on significativa del rendimiento entre el per\'iodo de entrenamiento y prueba:

\begin{table}[htbp]
\centering
\caption{Degradaci\'on del Sharpe Ratio: Entrenamiento vs Prueba}
\label{tab:in_out_sample}
\small
\begin{tabular}{@{}lrrr@{}}
\toprule
\textbf{Modelo} & \textbf{Sharpe Train} & \textbf{Sharpe Test} & \textbf{Degradaci\'on} \\
\midrule
Ridge & 1.12 & 0.93 & -17\% \\
CNN\_LSTM & 0.98 & 0.74 & -24\% \\
RandomForest & 1.45 & 0.39 & -73\% \\
XGBoost & 1.89 & 0.16 & -92\% \\
\bottomrule
\end{tabular}
\vspace{0.2cm}
\small\textit{
\small
Nota: Degradaci\'on = (Sharpe Test - Sharpe Train) / Sharpe Train. Mayor degradaci\'on sugiere mayor overfitting.}
\end{table}

Los resultados revelan:

\begin{itemize}
    \item \textbf{Ridge}: M\'inima degradaci\'on (-17\%), sugiriendo buen ajuste de regularizaci\'on.
    \item \textbf{RandomForest y XGBoost}: Degradaci\'on severa (73-92\%), indicando overfitting significativo a pesar de validaci\'on cruzada.
\end{itemize}

\subsubsection{Correcci\'on por Multiple Testing}

Con 23 modelos evaluados, el riesgo de falsos positivos por comparaciones m\'ultiples es sustancial. Aplicamos la correcci\'on de Bonferroni:

\begin{equation}
\alpha_{\text{Bonferroni}} = \frac{\alpha}{m} = \frac{0.05}{23} \approx 0.0022
\end{equation}

Bajo este criterio estricto, solo Ridge mantiene significancia estad\'istica (p-valor DM = 0.0044).

\subsubsection{False Discovery Rate (FDR)}

Alternativamente, el procedimiento de \cite{benjamini1995controlling} controla la proporci\'on esperada de falsos positivos entre los rechazos:

\begin{enumerate}
    \item Ordenar los p-valores: $p_{(1)} \leq p_{(2)} \leq \ldots \leq p_{(m)}$
    \item Encontrar el mayor $k$ tal que $p_{(k)} \leq \frac{k}{m} \cdot q$
    \item Rechazar H$_0$ para los $k$ menores p-valores
\end{enumerate}

Con $q = 0.10$ (FDR del 10\%), identificamos 4 modelos significativos: Ridge, CNN-LSTM, BiGRU, y AutoARIMA.

\subsection{Resumen de Validaci\'on Estad\'istica}
\label{subsec:validation_summary}

\begin{table}[htbp]
\centering
\caption{Resumen de Tests Estad\'isticos por Modelo}
\label{tab:validation_summary}
\small
\begin{tabular}{@{}lccccc@{}}
\toprule
\textbf{Modelo} & \textbf{P(S$>$0)} & \textbf{DM Test} & \textbf{CW Test} & \textbf{MCS} & \textbf{FDR} \\
\midrule
Ridge & 99.8\% & \cmark & \cmark & \cmark & \cmark \\
CNN\_LSTM & 99.1\% & \cmark & \cmark & \cmark & \cmark \\
BiGRU & 98.4\% & \cmark & \cmark & \cmark & \cmark \\
GradientBoosting & 96.8\% & -- & -- & \cmark & -- \\
\bottomrule
\end{tabular}
\vspace{0.2cm}
\small\textit{
\small
\item \cmark = pasa el test al nivel de significancia est\'andar (5\% o 10\% seg\'un el test). -- = no pasa.}
\end{table}

\subsection{Conclusiones de la Validaci\'on}
\label{subsec:validation_conclusions}

La validaci\'on estad\'istica proporciona evidencia mixta:

\begin{enumerate}
    \item \textbf{Evidencia positiva}: Ridge, CNN-LSTM y BiGRU demuestran capacidad predictiva estad\'isticamente significativa bajo m\'ultiples tests.

    \item \textbf{Cautela necesaria}: Los intervalos de confianza del Sharpe Ratio se solapan sustancialmente con el benchmark, sugiriendo que la superioridad no es robusta.

    \item \textbf{Overfitting parcial}: Modelos ensemble (RandomForest, XGBoost) muestran degradaci\'on severa entre entrenamiento y prueba.

    \item \textbf{Modelos regularizados}: Ridge y redes con dropout (CNN-LSTM, BiGRU) mantienen mejor generalizaci\'on.
\end{enumerate}

Estos hallazgos son consistentes con la literatura que sugiere que modelos m\'as simples y regularizados tienden a generalizar mejor en predicci\'on financiera \citep{gu2020}.
